% Original BEAV documentation converted to LaTeX source by
%
% Gary Bartlett (garyb@abekrd.co.uk)		22-Nov-1991

\documentstyle[a4,12pt,titlepage,twoside,fancyheadings]{article}
\addtolength{\oddsidemargin} {-6pt} 
\addtolength{\evensidemargin}{-14pt}
\setlength  {\topmargin}     {0pt}
\setlength  {\textheight}    {44\baselineskip}
\addtolength{\textheight}    {\topskip}

\pagestyle{fancy}
\lhead{}
\chead{BEAV User Manual}
\rhead{}
\lfoot{}
\cfoot{\thepage}
\rfoot{}
\addtolength{\headheight}{3pt}

% Tabbing for command descriptions
\newcommand{\cmd}[3]{
	\begin{tabbing}
	67890\=12345678901234567890123\=45678901234567890\= \kill
	     \> #1                    \> #2              \> #3
	\end{tabbing}
}

% Tabbing for key descriptions
\newcommand{\key}[3]{
	\begin{tabbing}
	67890\=123456789012345\=6789012345678901234567890\= \kill
	     \> #1            \> #2                      \> #3
	\end{tabbing}
}

\newcommand{\BEAV}{{\bf BEAV}}
\newcommand{\UNIX}{{\sc unix}}

\title{THE \\ BEAV \\ Binary Editor And Viewer}

\author{Manual Copyright 1991 \\[1cm] By \\ Peter Reilley \\
	19 Heritage Cir. \\ Hudson, N.H. 03051 \\ \tt pvr@wang.com \\[1cm]
	BEAV source and executable can be freely \\
	distributed for non-commercial purposes.}

\date{Version 1.32 \\ November 8, 1991
	\vfill\newpage\thispagestyle{empty}}

\begin{document}

\maketitle

\setcounter{page}{1}

\section{Introduction}

      \BEAV\ is an editor that brings the features of a powerful
full screen editor to the editing of binary files.   It is the
only editor that I know of that does that.

      When you need to edit a non-text file you generally have
two choices; a text editor or a file zap type editor.   Each
choice has significant disadvantages.

      Text editors expect the file to be formatted in a certain
way.   At a minimum they expect that all lines be terminated by a
carriage return or line feed and be limited in length.   There is
no line length limit with \BEAV.   Most text editors get confused
by bytes that are outside of the normal range (20 to 7E HEX).
In \BEAV\ no special characters such as carriage return or line
feed affect the display aside from producing their numeric value.
\BEAV\ can edit any file no matter the format.

      The other choice is to use a file zap type editor which can
edit a binary file without difficulty.  These editors are often
very limited in their features and capabilities.   Most file zap
programs can edit a file only in HEX or ASCII.   They generally
operate on a sector basis and because of this they cannot insert
or delete data in the middle of the file.

      All these limits are eliminated in \BEAV.   You can edit a
file in HEX, ASCII, EBCDIC, OCTAL, DECIMAL, and BINARY.   You can
search or search and replace in any of these modes.   Data can be
displayed in BYTE, WORD, or DOUBLE WORD formats.   While
displaying WORDS or DOUBLE WORDS the data can be displayed in
INTEL's or MOTOROLA's byte swap format.   Data of any length can
be inserted at any point in the file.   The source of this data
can be the keyboard, another buffer, of a file.   Any data that
is being displayed can be sent to a printer in the displayed
format.   Files that are bigger than memory can be handled.

      Some users may recognize the similarity to the EMACS text
editor that was written by Richard Stallman at MIT.   This is not
a coincidence.   I attempted to keep as much of the user
interface and functionality as possible given the very different
tasks of EMACS and \BEAV.

\subsection{Quick Start}

      This is a brief description of the minimal set of commands
that are necessary to start using \BEAV\ effectively.

      The file-visit command, Ctl-X Ctl-V, can be used to read a
file in for editing.   The file can also be read in from the
command line; {\tt beav $<$edit\_file$>$}.

      These commands can be used to navigate around the window.

	\cmd{move-back-char}{Ctl-B}{moves left}
	\cmd{move-back-line}{Ctl-P}{moves up}
	\cmd{move-forw-char}{Ctl-F}{moves right}
	\cmd{move-forw-line}{Ctl-N}{moves down}

      The move-to-byte command, Ctl-X G, will prompt you for a
byte position to move to.

      These commands will insert a zero byte at the cursor
position or delete the byte at that position.

	\cmd{insert-unit}{Ctl-X I}{}
	\cmd{delete-forw-unit}{Esc D}{}

      The file-save command, Ctl-X Ctl-S, will save the data to
the file if a change has been made.

      The help command, Esc ?, will display a list of all
commands and their current key bindings.

      The abort-cmd command, Ctl-G, will abort any command that
is in operation.

      The quit-no-save command, Ctl-X Ctl-C, will exit \BEAV.   If
there is any data that has not been saved you will be warned.


\section{Overview}

\subsection{Terms and Definitions}

      Throughout this manual certain terms will be used to
describe the operation and structure of \BEAV.

      The data that \BEAV\ is editing is held in a buffer that is
structured as a byte stream.   There are many commands that
affect the way that this byte stream is displayed on to the
screen.   Whatever display mode is chosen the data in the buffer
is not effected, only the presentation.

      One such choice is to display the data as bytes, words, or
double words.   That is; 8 bit values, 16 bit values, or 32 bit
values.   Whatever choice is made the value of the selected size
will be displayed.   These values are referred to as units in
this manual.   Thus the `delete-forw-unit' command deletes the
unit under the cursor.   If 32 bit units are being displayed then
4 bytes will be deleted.

      Many commands in \BEAV\ start by pressing the `ESCAPE' key.
When this manual refers to a command that requires that the
`ESCAPE' key be pressed it will be abbreviated with `Esc'.
Another frequently used key stroke to start commands is `CONTROL
X'.   This in done by pressing the `CONTROL' key then pressing
the `X' key at the same time.   In this manual this will be
abbreviated by `Ctl-X'.   Many commands contain other control
characters and these will be abbreviates similarly.   Thus the
`insert-unit' command will be listed as `Ctl-X I'.   This will be
entered be pressing the CONTROL key and while holding it hitting
the `X' key, release the CONTROL key then hit the `I' key.

\subsection{The Screen Format}

   \BEAV\ presents information to the user in a number of areas.
The first is the window.   There will be at least one window
displayed on the screen at all times.   The window consists of
two areas.   The first is the display area.   This is where the
data that is in the buffer is displayed.   Most of the time the
cursor will be in this area, as this is where most editing is
done.  Each line in the display area will start with a number
that indicates the offset into the buffer for this line of data.
At the bottom of the display area for each window is the status
line.

     The status line presents the user with a number of specific
pieces of information.   The first is the program name which is
``BEAV''.   Next there are some flags indicating the status of this
particular buffer.

     The first flag indicates if the buffer has been damaged by a
memory allocation failure.   If there is a `?' then the buffer is
bad, otherwise there will be a space.   Bad buffers are set to
read only.

     The next flag indicates if the buffer has been changed and
will need to be written out if the changes are to be saved.
This flag is a `$\ast$' if the buffer has been changed, otherwise
there will be a space.

     The next flag indicates insert/overstrike mode; `I' for
insert, `O' for overstrike.  This is only meaningful in ascii or
ebcdic mode.

     Then the buffer name followed by the file name.   A window
can be in read only, read/write, or read/write with size lock.
These  will be displayed as [RO], [RW], or [WL] respectively.
If auto-update is enabled then [AU] will be displayed.

     Next the cursor position in bytes and the character position
within the unit.   The next piece of information gives the format
that the data is being displayed in; HEX, OCTAL, DECIMAL, BINARY,
ASCII, EBCDIC.   If a data mode is selected then the size in
bytes (1, 2, 4) is displayed.  If the data is being  displayed
in Intel mode then the swapped flag is displayed, `S'.   Lastly
the byte shift is displayed; 0 only for 8 bit data, 0 or 1 for 16
bit data, or 0, 1, 2 or 3 for 32 bit data.

     There can be multiple windows on the screen at the same time
but each window will have a status line at the bottom.

     The very bottom line on the screen is the prompt line.
This is where you enter data that \BEAV\ requests.   If you want to
edit a new file you would type `Ctl-X Ctl-V', \BEAV\ would respond
with ``Visit file:'' on the prompt line.   The cursor would be
positioned after the prompt.   You would then enter the name of
the file that you wished to edit.

     If you entered the command by mistake, you can abort the
operation by typing a `Ctl-G'.   `Control G' is a universal abort
command and can be used anywhere.   If you want to perform a
search you will enter the search string on this line.   When you
have entered the information that \BEAV\ has requested hit `Return'
and the cursor will return to it's original position in the
window display area.   The prompt line is also where error
messages are displayed.

\subsection{Display Modes}

      The data in the buffer can be displayed in a number of
formats.   First there is the display mode.   This can be either;
HEXADECIMAL, DECIMAL, OCTAL, BINARY, ASCII, or EBCDIC.

      If ASCII or EBCDIC mode is chosen then each byte in the
buffer will be converted to it's ASCII or EBCDIC character and
displayed.   Bytes that are outside of the standard character set
will be displayed as a dot.   Each line will be 64 characters
long.   The byte value for ``carriage return'' and ``line feed'' will
be displayed as a dot as will any other non-printable characters.

      Within HEXADECIMAL, DECIMAL, OCTAL, or BINARY format the
data can be displayed in 8, 16 or 32 bit values.   If 16 or 32 it
values are displayed the user can choose to view these values in
either the Intel format or the Motorola format.   If Intel format
is chosen then the first byte in each unit is the least
significant byte when the value is calculated.   Thus in hex 32
bit mode a byte string of ``32 7A 9F 10'' would be displayed as
``109F7A32''.   If Motorola format is chosen this value would be
displayed as ``327A9F10''.

      There is another display format choice that affects the 16
or 32 bit formats.   This is called shift.   The shift can be 0
or 1 for 16 bit modes, or 0, 1, 2, 3 for 32 bit modes.   Shift
moves the zero point reference for the buffer up by the selected
value.   The default is zero shift.   If a buffer is displaying
the following 32 bit hex data;
	\begin{center}
           ``12345678 2F4792AF 673DFEA1 88551199''
	\end{center}
with the shift at 0.   Changing shift to 1 will produce;
	\begin{center}
           ``3456782F 4792AF67 3DFEA188 55119955''
	\end{center}
The data has been slid down toward the beginning of the buffer by
one byte.   This has not changed the data in the buffer at all,
it has only affected the way that the data is presented on the
screen.   This is useful when looking at WORD or DOUBLE WORD data
that is not aligned on two or four byte boundaries.

      When \BEAV\ is first started or a new window is opened the
default format is HEXADECIMAL BYTES.

\subsection{Commands}

      Commands are the means that the user controls the operation
of \BEAV.   A command can be given by using one of two methods.
The first is to use the key binding.

     A command can have one or more associated key bindings.   It
can also have no key binding.   There are a set of default key
bindings that \BEAV\ comes configured with.   The current set of
key bindings can be seen by using the `help' command.   The
`help' command is `Esc ?' or Function Key 1 `F1' on the IBM PC.
The help buffer can be scrolled by using the up and down arrow
keys.   A printed copy may be obtained by going to the bottom of
the help buffer using the `move-to-end' command (`Esc $>$' or the
`End' key).   Then issue the `print-mark-to-cursor' command (`Esc
P' or `Ctl-Print') and enter `PRN' or a file name when prompted
with ``Print to:''.   This will output the entire help buffer to a
printer connected to the parallel interface or the specified
file.

     The second method of issuing a command is to use the
`extended-com\-mand' command (Esc X).   You are prompted for a
command name that you want to execute.   This method is useful
for executing commands that have no key binding.   Commands that
can take a long time to complete can be stopped by pressing Ctl-
G.

\subsection{Buffers}

      Buffers are the in memory storage for all data editing and
viewing.   Each buffer has a name that appears in the mode line.
Buffers generally have a file name that is associated with them.
The file name also appears in the mode line.   The buffer name
and the file name are independent but the buffer name defaults to
the file name.

      The buffer name is used to refer to a specific buffer.
The `change-buffer' (`Ctl-X B') command will prompt you for a
buffer name.   After you enter a buffer name that buffer will be
displayed in the current window.   If there is no such buffer,
one will be created and displayed (it will be empty).

      When \BEAV\ is run with a file name as a command line
parameter, the file is read into a new buffer.  The buffer name
will be made the same as the file name.   The file name is only
used when the buffer is saved.   If the file name is changed
using the `buffer-set-file-name' (`Ctl-X Ctl-F') command then
when the buffer is saved it will be saved to the new file.

      Buffers are dynamically allocated.   They grow or shrink as
the size of the data they contain changes.   The buffer size can
be frozen using the `buffer-size-lock' (`Ctl-X Ctl-L') command.
This prevents inserting or deleting data from the buffer but data
can be modified.

      Buffers continue to exist even if they are not being
displayed.   Buffers are saved in a linked list that can be
stepped through using the `change-to-next-buffer' (`Esc +') or
`change-to-prev-buffer' (`Esc -') commands.   If the `change-to-
next-buffer' command is given then the next buffer in the list is
displayed in the current window.

\subsection{Files}

      Files are the means of storing data on disk.   Files or
segments of files can be loaded into \BEAV\ for editing or viewing.
The file on the disk has not been changed until \BEAV\ does a save
to that file.   When a save to a file is performed the original
file contents in saved in a ``{\tt .bak}'' file.

\subsection{Key Binding}

      All commands in \BEAV\ have a command name and a default key
binding.   The bindings may be changed to suit personal
preference by using the `bind-to-key' (`Esc K') command.   The
current binding of a key sequence can be displayed by using the
`binding-for-key' (`Ctl-X ?') command.

      Key bindings can be loaded automatically from a file named
``{\tt .beavrc}'' under \UNIX\ or ``{\tt beav.rc}'' under MSDOS.   I will refer to
this file as the {\tt beavrc} file in this document.   This file must
be in a directory specified by the PATH environment variable.
When \BEAV\ is started this file is read and the bindings are
loaded.   This file is a simple text file and can be edited to
make changes.

      The {\tt beavrc} file contains three columns.   The first is the
function name to be bound.   The second is the key sequence that
is bound to that function.   The third is a number that defines
the key code.   When the {\tt beavrc} file is read, in only the first
and third columns are used.

      The simplest way to create a valid {\tt beavrc} file is to first
bind the key codes to the desired functions using the `bind-to-
key' command (Esc K).   Next display the current bindings in a
window using the `help' command (Esc ?).   Now save that buffer
to a file with the `file-write' command (Ctl-X Ctl-W).   You will
be prompted for a file name.   The file written is a valid {\tt beavrc}
file.   You may want to edit the {\tt beavrc} file to remove the
commented lines, ie, those beginning with \#, to make it load
faster.

      On \UNIX\ systems there can be multiple {\tt .beavrc} files, one
for each terminal type.   This is accomplished by appending the
{\tt TERM} variable to {\tt .beavrc}.   Thus if you use both a vt100 and a
wyse60 terminal, you could have a {\tt .beavrc.wy60} and a
{\tt .beavrc.vt100} bindings file.   When your {\tt TERM} variable is
{\tt TERM=wy60} then the {\tt .beavrc.wy60} bindings file will be used.   If
there is no {\tt .beavrc.wy60} file then the {\tt .beavrc} file will be used.

\subsection{Configuration}

      When the MSDOS version of \BEAV\ run it detects whether the
system is an IBM PC or a clone.   If a PC is detected then a set
of key bindings that use the 10 function keys and the relevant
keypad keys are loaded.   If the system is not recognized then
only the standard bindings are loaded.

      If a PC is detected the screen is managed by making BIOS
level calls that enhance performance.   Otherwise, the screen is
controlled by making ANSI compatible calls to the operating
system.   This is much slower but is not sensitive to hardware
configuration.   This requires that non-standard MSDOS systems
support ANSI display controls.   The following ANSI escape
sequences are used;

      \begin{center}\begin{tabular}{ll}
      Position cursor          & ESC [ $<$row$>$ ; $<$column$>$ H	\\

      Erase to end of line     & ESC [ 0 K	\\

      Erase to end of page     & ESC [ 0 J	\\

      Normal video             & ESC [ 0 m	\\

      Reverse video            & ESC [ 7 m
      \end{tabular}\end{center}

      On \UNIX\ systems the termcap library is used.   This
requires that the {\tt TERM} environment variable be set correctly.


\section{Command Description}

      This section describes all \BEAV\ commands as to their
function and any side effects that they may have.   The first
line of each description of a command will begin with the default
key binding then the command name and follows with the key
binding for a PC.

\subsection{Help}

      This command returns information that will aid in the use
of \BEAV.

	\key{Esc ?}{help}{F1}

      A new window is opened by splitting the current window then
all current key bindings are displayed.   This buffer is larger
than can be shown at one time and must be scrolled up and down to
see all entries.    All commands that do not alter data can be
used to move around and search the help buffer.   To leave the
help buffer use the `window-delete' command (`Ctl-X 0').

      Lines that begin with the comment character (\#) are
commands that have the default bindings.   Lines without the
comment character have been bound by the user or by processing a
{\tt beavrc} file.   The comment character helps in creating a {\tt beavrc}
file by identifying only those bindings that have changed.

\subsection{Cursor Movement}

      There are a number of commands that move the cursor around
the current window.   If the cursor bumps the top or the bottom
of the window the position of the window will be adjusted so as
to keep the cursor within the bounds.   When the window is moved
in response to cursor movements it is moved by about one third of
the window size.   This improves performance by reducing the
number of window moves.

	\key{Ctl-P}{move-back-line}{North (up arrow)}
	\key{Ctl-N}{move-forw-line}{South (down arrow)}

      These commands move the cursor up one line or down one
line.   If the cursor is on the top line in the buffer and a
`move-back-line' command is given the cursor will jump to the
beginning of the first unit in the buffer.   If the cursor is on
the last line of the buffer and a `move-forw-line' is given the
cursor will move to the last unit in the buffer.

	\key{Ctl-F}{move-forw-char}{East (right arrow)}
	\key{Ctl-B}{move-back-char}{West (left arrow)}

      These commands move the cursor forward or backward in the
current line.   If the cursor is at the first character in the
first unit of the line and the `move-back-char' command is given
then the cursor will wrap to the last character of the previous
line.   If the cursor is at the last character of the last unit
in the current line then it will wrap to the first character of
the next line.

	\key{Esc F}{move-forw-unit}{Ctl-East}
	\key{Esc B}{move-back-unit}{Ctl-West}

      These commands are similar to the above set but they move
the cursor by units rather than characters.   The command `move-
forw-unit' will position the cursor to the first character of the
next unit.   The command `move-back-unit' will move the cursor to
the first character of the previous unit.

	\key{Ctl-V}{move-forw-page}{PageUp}
	\key{Esc V}{move-back-page}{PageDown}

      These commands move the move the data in the window by the
number of lines in the window less one.   The cursor will stay in
the same position relative to the window as the data is moved.

	\key{Esc $<$}{move-to-beginning}{Home}
	\key{Esc $>$}{move-to-end}{End}

      Move the cursor to the beginning or the end of the buffer.

	\key{Ctl-X G}{move-to-byte}{F9}

      Prompt for a byte offset, then go to that position in the
current buffer.

	\key{Ctl-X Ctl-N}{move-window-down}{Ctl-Z}
	\key{Ctl-X Ctl-P}{move-window-up}{Esc Z}

      Move the buffer in the window up or down by one line.
This does not effect the cursor until it hits the top or bottom
of the window.

	\key{Esc .}{mark-set}{F2}

      Set the mark position to the current cursor position.   The
mark position is remembered even for nonactive windows and
buffers.

	\key{Ctl-X Ctl-X}{swap-cursor-and-mark}{}

      The position of the cursor and the position of the mark are
swapped.

	\key{Esc L}{window-link}{}

      This command causes all windows that are displaying the
contents of the current buffer to have the same cursor position.
Thus if one window is scrolled then all other windows that
display that buffer are scrolled so as to keep the cursor in the
window.

	\key{Ctl-X =}{show-position}{}

      The current position of the cursor and the mark are
displayed.   The buffer size, file size and file name are also
shown.

\subsection{Buffer Management}

      Buffers store all data that is being edited.   The only
limit to the number of buffers is the size of available memory.
If a file is loaded into a buffer and there is insufficient
memory to hold the entire file, then it will be loaded until
memory is exhausted.   The buffer will then be set to read only
mode.

	\key{Ctl-X Ctl-B}{buffers-display}{Ctl-F1}

      A new window is opened and a list of all buffers in \BEAV\ is
displayed.   The list contains the buffer name, the file name (if
it exists), the buffer size, and a state flag.   If the list is
too large for the window, you can scroll the list.   If you
position the cursor on a line describing a buffer and give the
set-mark (Esc .) command you will be prompted to; Go to, Kill, or
Save that buffer.

	\key{Ctl-X B}{change-buffer}{Ctl-F2}

      This command prompts you for a buffer name.   If you enter
the name of an existing buffer, that buffer will be displayed in
the current window.   If the name does not match an existing
buffer, a new buffer will be created and displayed.   This buffer
will be empty and will have no file name.

	\key{Esc +}{change-to-next-buffer}{Ctl-F4}
	\key{Esc -}{change-to-prev-buffer}{Ctl-F5}

      The next or the previous buffer in the buffer list is
displayed in the current window.   This does not effect buffers
that are displayed in other windows.

	\key{Esc G}{move-to-buffer-split}{}

      Prompt for a buffer name.   Then split the current window
and display that buffer, creating it if necessary.

	\key{Esc Ctl-N}{buffer-set-name}{Esc Ctl-N}

      The current buffer name is changed to the name that you
enter.   This does not effect the file name.

	\key{Ctl-X Ctl-F}{buffer-set-file-name}{Ctl-F7}

      The file name of the current buffer is changed to the name
that you enter.   This does not affect the buffer name.

	\key{Ctl-X K}{kill-buffer}{Ctl-F3}

      This command prompts you for a buffer name.   This buffer
is then deleted.   If the buffer is currently being displayed you
are prompted for conformation.   If the buffer has been changed
you are again prompted for conformation.

	\key{Ctl-X Ctl-L}{buffer-size-lock}{}

      The buffer size is prevented from being changed.   Data can
be edited but only by changing the existing data.   If a buffer
is copied into a size-locked buffer the operation well be
successful but will overwrite existing data.   This command
toggles between locked and unlocked.

	\key{Esc Y}{yank-buffer}{Ctl-F6}

      Data from one buffer is inserted into the current buffer at
the cursor position.   You are prompted for the name of the
buffer to copy from.

	\key{Esc O}{save-mark-to-cursor}{}

      Prompt for a buffer name.   Create a new buffer with that
name and write the data from the mark to the cursor into that
buffer.

	\key{Esc Ctl-W}{show-save-buf}{}

      Displays the contents of the save buffer in a new window.
The save buffer cannot be edited.   This command can be used to
see the data that the yank (Ctl-Y) command will insert into the
current buffer.

\subsection{File Management}

      These commands control \BEAV's access to files.   Files are
loaded into buffers or are written from buffers.   Commands that
prompt for a file name also accept range parameters.   Range
parameters are always given in the numeric base of the current
window.   Thus if you are displaying data in decimal format then
the range parameters must be entered in decimal.

      The size of a file read or write can be limited by
specifying a range.   The range parameter specifies the offset
into the file, not the buffer.   Range parameters can be
specified in these forms;

      \begin{tabular}{lll}
      $<$file name$>$ & $<$start address$>$ &	\\

      $<$file name$>$ & $<$start address$>$ & $<$end address$>$	\\

      $<$file name$>$ & $<$start address$>$ & +$<$length$>$
      \end{tabular}

      The first form causes the read or write to begin from the
$<$start address$>$ value until the end of the buffer on write or the
end of the file on read.

\begin{sloppypar}
      The second form reads or writes from $<$start address$>$ until
$<$end~address$>$ non-inclusive.
\end{sloppypar}

      The third form reads or writes from $<$start address$>$ for
$<$length$>$ bytes.

      Thus, if the command `file-read' is given and you enter at
the prompt;   {\tt main.obj 1000 +100}.   If the current display mode
is hex, the file ``{\tt main.obj}'' will be read from hex byte address
1000 to 10FF into the buffer.

	\key{Ctl-X Ctl-R}{file-read}{Sh-F2}

      Prompt for a file name and read that file into the current
buffer.   This overwrites the current contents of the buffer.
The buffer name is not changed but the buffer file name is set to
the new file name.

	\key{Ctl-X Ctl-S}{file-save}{Sh-F3}

      Write the current buffer out to the file if it has been
changed.   If the buffer has not been changed then do nothing.

	\key{Ctl-X V}{file-view}{}

      Prompt for a file name and read file into a new buffer and
display in current window.   Set to read-only mode.

	\key{Ctl-X Ctl-V}{file-visit}{Sh-F4}

      Prompt for a file name.   If the buffer already exists then
display it in the current window.   Otherwise, read file into a
new buffer and display in current window.   If there is no such
file then create it.

	\key{Esc U}{file-visit-split}{}

      Same as above but split current window and display new
buffer.   This displays the new buffer as well as the old buffer.

	\key{Ctl-X Ctl-W}{file-write}{Sh-F5}

      Prompt for a file name, then write the current buffer to
that file.

	\key{Ctl-X Tab}{insert-file}{Sh-F8}

      Prompt for a file name and insert that file into the
current buffer at the cursor position.

	\key{Ctl-X Return}{save-all-buffers}{Sh-F6}

      Write all buffers that have been changed to their
respective files.

\subsection{Window Management}

      \BEAV\ presents information to the user in one or more
windows.   Each window is a view into a buffer where the data is
actually stored.   The window controls how this data is formatted
for display to the user.   Data can be displayed as HEX bytes,
OCTAL bytes, ASCII characters, plus many others.   The display
format is associated with the window.   Thus if a new buffer is
displayed in the current window that new data will be displayed
in the current windows format.

      The only limit to the number of windows is the screen size.
A window can be no smaller than two lines.   This along with the
mode line associated with each window limits to eight the number
of windows on an IBM PC 25 line screen.

      Any window can view any buffer including having many
windows on the same buffer.   For example, two windows can
display the same buffer but present the data in two different
modes.   One window could display HEX bytes and the other could
display ASCII characters.

	\key{Ctl-P}{change-window-back}{Ctl-PageUp}
	\key{Ctl-N}{change-window-forw}{Ctl-PageDown}

      These commands move the cursor to the next or previous
window on the screen, making it the current window.

	\key{Ctl-X Z}{window-enlarge}{}
	\key{Ctl-X Ctl-Z}{window-shrink}{}

      Enlarge or shrink the current window size by one line.

	\key{Esc !}{window-reposition}{}

      Move window position so that the cursor is centered in the
window.   The cursor position in the buffer does not change.

	\key{Ctl-X 2}{window-split}{}

      Split the current window into equal parts.   Both haves
have the same display mode and view the save buffer.

	\key{Ctl-X 1}{window-single}{}

      Expand the current window to fill the entire display, all
other windows are removed.   Make the current window the only
window displayed.   This has no effect on the underlying buffers
except that they may not be displayed any more.

	\key{Ctl-X 0}{window-delete}{}

      Delete the current window and expand the upper window to
fill the space.   This has no effect on the underlying buffer
except that it may not be displayed any more.

	\key{Esc Ctl-A}{display-ascii}{}
	\key{Esc Ctl-E}{display-ebcdic}{}
	\key{Esc Ctl-H}{display-hex}{}
	\key{Esc Ctl-O}{display-octal}{}
	\key{Esc Ctl-D}{display-decimal}{}
	\key{Esc Ctl-B}{display-binary}{}

      These commands set the window display mode.   Text buffers
can be displayed as ASCII or EBCDIC characters.   Buffers that
are not human readable can also be displayed in hexadecimal,
octal, decimal, or binary format.

	\key{Esc 1}{display-bytes}{}
	\key{Esc 2}{display-words}{}
	\key{Esc 4}{display-double-words}{}

      As a further option on the non-text display modes, data can
be displayed as 8, 16, or 32 bit values.

	\key{Ctl-E}{display-swap-order}{}

      When data is displayed as words or double words the order
of significance can be changed.   In Intel microprocessors the
least significant byte of a word is stored at the lowest address.
Thus if the word 5892 (HEX) were stored at memory address 10,
then 92 (HEX) would be stored at address 10 and 58 (HEX) would be
stored at address 11.   In Motorola microprocessors the reverse
is true.

      This command toggles between the Intel and Motorola schemes
of assembling bytes into words and double words.   This command
has no effect on byte display or on the text display modes.   The
data in the buffer is not changed.

	\key{Ctl-A}{display-byte-shift}{}

      This command changes the offset from the beginning of the
buffer used to assemble words and double words.   The default
shift is 0.   For example, a double word at address 10 is made up
of the bytes at address 10, 11, 12, and 13.   With a shift of 1
that double word would be made of bytes 11, 12, 13, and 14.
With a shift of 2 then bytes 12, 13, 14, and 15 would be used.
The maximum shift in word display mode is one and the maximum
shift in double word mode is three.

      The buffer is in effect shifted toward the beginning of the
buffer with 1, 2, or 3 bytes becoming not visible.   These bytes
are not lost, they become visible when the shift is set to zero.
This command cycles through all possible shift values.   There is
no effect in any byte display mode or any text display mode.

\subsection{Inserting and deleting}

      These commands are the core of \BEAV.   These commands allow
the buffer to be edited in a similar fashion to a text editor.
\BEAV\ has an insert mode much the same as text editors but it only
works when displaying data in one of the text modes, either ASCII
or EBCDIC.   In other modes it doesn't make any sense to insert
characters as they are typed when there is more than one
characters per unit.   In the data modes there is a command that
inserts a unit of zeros into the buffer.   Similarly the delete
commands always delete a unit rather than a character.   In a
text mode the delete commands work as in a text editor because a
unit is a character.

	\key{Ctl-X I}{insert-unit}{}

      Insert a zero at the cursor position.   The rest of the
data moves down one place.   Thus, if double words are being
displayed, four bytes are inserted before the cursor position.
These bytes are initialized to zero.   This command works in all
display modes.

	\key{Ctl-X Ctl-A}{insert-toggle}{Insert}

      In either of the two text modes this command toggles
between insert mode and overwrite mode.    In insert mode each
character that is typed is inserted in front of the cursor and
the rest of the buffer is moved down.   In overwrite mode the
typed characters replace the character that is at the cursor.
This command has no effect in a non-text display mode.

	\key{Ctl-Q}{insert-literally}{Esc Q}

      This command sets a special temporary mode where the next
typed character is inserted in the buffer no matter what the
character is.   This allows control codes to be inserted in the
buffer when in a text display mode.   Alternatively the same byte
could be inserted into the buffer by using one of the data
display modes.   It night be faster to use this command on some
occasions.

	\key{Ctl-T}{unit-twiddle}{}

      The unit at the cursor is swapped with the previous unit.

	\key{Rubout}{delete-back-char}{Backspace}

      This command deletes the character before the cursor and
pulls the rest of the buffer back.   The cursor remains on the
same character as it moves back.   It only works in the text and
byte display modes.

	\key{Ctl-D}{delete-forw-char}{Delete}

      The character at the cursor is deleted and the buffer is
pulled back.   The cursor remains at the same position.   It only
works in the text and byte display modes.

	\key{Esc Rubout}{delete-back-unit}{Esc Ctl-K}

      This command deletes the unit before the cursor and pulls
the rest of the buffer back.   The cursor remains on the same
unit as it moves back.

	\key{Esc D}{delete-forw-unit}{}

      The unit at the cursor is deleted and the buffer is pulled
back.   The cursor remains at the same position.

	\key{Esc W}{copy-mark-to-cursor}{F7}

      The area in the buffer from the mark to the current cursor
position is copied into the kill buffer.   If the mark is not set
before this command is given an error is reported.

	\key{Ctl-W}{delete-mark-to-cursor}{F8}

      The area in the buffer from the mark to the current cursor
position is deleted and placed into the kill buffer.   If the
mark is not set before this command is given an error is
reported.

	\key{Ctl-Y}{yank}{F6}

      The contents of the kill buffer is inserted into the buffer
at the cursor position.   The kill buffer itself is not changed.

\subsection{Search and Replace Commands}

      \BEAV\ has very powerful search and replace commands.   The
search and replace string can be entered in any of the display
modes.   The search and replace strings can each be up to 256
bytes long.   The display mode can be changed at any time while
entering the string.   Wild cards can be placed down to the bit
level on both the search and replace strings.   The wild card
character, `?', will match any value that it is compared with.

      When a wild card is placed in the replace string it leaves
the destination data unchanged at that position.   Thus, if the
destination contains the ASCII string `41 42 43 44' and the
replace string contains `66 67 ?? 69' the result would be `66 67
43 69'.

      Wild cards can be placed in any position that makes sense.
If you want to use wild cards in an ASCII string then you must
switch to another mode to enter them.   You can then switch back
to ASCII mode.   In this case a `?' will appear in the position
where the wild card has been placed but it appears exactly the
same as a standard question mark.   In fact if you type a `?'
over the wild card there will be no apparent change.   However,
the character will no longer be a wild card but a standard
question mark.   To see the true wild cards you must use a data
display mode.   In fact if the wild card has been set on the bit
level then you must go to binary display mode to see its actual
position.

      The commands to change the display mode in search and
replace are the same as for the window display mode.   The search
and replace strings can be scrolled back an forth and the cursor
moved using the same commands as for the window.   While
performing a replace command you can switch between the search
string and replace string by using the `move-back-page' or `move-
forw-page' commands.

	\key{Esc S}{search-forw}{F3}

      Prompts for a search string then searches from the current
cursor position for the first match.   The cursor is positioned
at the first unit of the match.

	\key{Esc R}{search-back}{}

      This command is the same as the previous one except that it
searches backward.

	\key{Esc T}{search-again}{F4}

      This command repeats the previous search command, forward
or backward.   The cursor is first moved one byte in the
appropriate direction before the search is repeated.

	\key{Esc \%}{replace}{F5}

      Prompt for search string.   After entering the search
string hit return and you will be prompted for the replace
string.   After entering the replace string hit return.   \BEAV
will then search for the first match with the search string.   If
a match is found you will be prompted with `(R)eplace, (S)kip,
(A)ll, (O)ne, (Q)uit'

      If you type a `R' the replace will be done at this location
and the search will continue.   If you type a `S' the replace
will not be done and search will continue.   If you type an `A'
the replace will be done and will be done at all future matches
without pausing for conformation.   If you type an `O' the
replace will be done at this location and the search will stop.
If you type a `Q' then the search will be terminated.

	\key{Ctl-R}{recall-srch-string}{}

      If you enter search or replace previously used strings can
be recalled with this command.

\subsection{Exiting BEAV}

      While using \BEAV\ individual buffers may be saved to disk
during the editing session.   When quitting \BEAV\ you must save
all buffers or delete all buffers.   There are two commands that
do this.

	\key{Ctl-C}{quit-no-save}{Sh-F10}

      If there are any unsaved buffers you will be prompted for
conformation before proceeding.   All buffers will be deleted
then you will return to DOS.

	\key{Ctl-X Ctl-E}{quit-save-all}{Sh-F9}

      All buffers are saved before exiting to DOS.

\subsection{Printing}

      The data that is being displayed in \BEAV\ can be printed or
sent to a file in the same format as displayed.   If the current
window is displaying octal words and a print command is given the
format of the print will be in the format of the window; that is,
octal words.

	\key{Esc P}{print-mark-to-cursor}{Ctl-Print}

      To use this command you must set the mark and the cursor to
define the region that you want printed.   If the mark is not set
it as assumed to be at the first unit.   After you enter the
command you will be prompted with `Print to:'.   You can enter a
file name or a device name to send the print image to.   If you
enter `PRN' most systems will print a hard copy.

      This is useful for getting a print out of the current key
bindings.   To do this give the `help' command `F1'.   Go to the
bottom of the help window using the `move-to-end' command `End',
the mark will be assumed to be at the beginning of the buffer.
Issue the `print-mark-to-cursor' command.   Enter `PRN' at the
prompt.   This should print the complete help buffer and will
reflect any changes that you have made to the key bindings.

\subsection{Keyboard Macros}

      \BEAV\ has the capability of recording key strokes as they
are entered and playing them back later.   This is useful for
repeating multi-keystroke operations.

	\key{Ctl-X (}{macro-start}{}

      Start recording key strokes.   There is no effect on the
operation of the key strokes.   Any previous recorded key strokes
are cleared.

	\key{Ctl-X )}{macro-end}{}

      Stop recording key strokes.   The key strokes are available
for play back.

	\key{Ctl-X E}{macro-execute}{}

      Play back the recorded key strokes.   The key strokes that
were recorded are played back as if they were typed at the
keyboard.

\subsection{Key Binding}

      \BEAV\ provides a user configurable interface.   The
interface is controlled by a set of key bindings.   This relates
the command that will be executed when a particular key stroke is
entered.   There are a set of default key bindings as described
in this manual.   These can be changed to reflect your
preferences.   When a change is made it is reflected in the help
screen.

	\key{Ctl-X ?}{binding-for-key}{Sh-F1}

      This command will tell you what function a certain key
sequence is bound to.   When this command is given you will be
prompted for a key stroke or key stroke sequence.   \BEAV\ will
report back with the function name.

	\key{Esc K}{bind-to-key}{}

      First you will prompted for a function name.   Enter the
name of the function that you wish to create a new binding for.
Function names are the names listed in this manual that are of
the form of `move-forw-unit' or `display-hex'.   After you enter
the name hit return.   You will be prompted for a key.   This can
be in the form of a single standard key such as `Z'.   Standard
key sequences can be entered such as `Ctl-X Z'
or `Esc Z'.   Special keys can be entered such as `F1' (function
key 1) or `Page Down'.   It is probably a good idea to not use
keys that are needed for editing.   If you bound `Z' to a
function then you would not be able to enter it as a keystroke
when using ASCII display mode.   You could still enter it using
the `insert-literally' command or doing it in one of the data
display modes but this would be more cumbersome.

	\key{Ctl-X L}{bindings-load}{}

      You are prompted for a file name that contains the key
binding that you wish to set.   This file is read in and the
appropriate bindings are set.   The text in the binding file
should be of the form;

	\key{$<$key name$>$}{$<$function name$>$}{$<$key code$>$}
      For example;
	\key{Ctl-X Ctl-P}{move-back-char}{0550}
	\key{F1}{move-forw-char}{04bb}
	\key{Ctl-A}{move-forw-unit}{0141}
	\key{Esc Ctl-T}{move-back-unit}{0354}

      The easiest way of producing a valid key binding file is to
set the desired bindings in \BEAV.   Next issue the `help' command
(ESC ?), then write the buffer out with the file-write command
(Ctl-X Ctl-W).   The file created will be a valid format for
loading and can be edited as desired.   It is the only reliable
way to get the $<$key code$>$ number.

\subsection{Special Functions}

      These are the commands that do not logically fit under one
of the previous headings

	\key{Gtl-G}{abort-cmd}{F10}

      This command aborts the current command.   It can even
abort a partially entered command.   Thus, if you have typed an
`Esc' as that start of a command you can type Ctl-G to return to
the normal command entry mode.

	\key{Esc A}{auto-save}{}

      \BEAV\ can be set to automatically save the current buffer
after a specified number of buffer editing commands are given.
This command first prompts for the number of operations before
the save is made.   If a zero is entered at the prompt, this
feature is disabled.   The default condition of this command is
disabled.

	\key{Esc C}{compare}{}

      This is a powerful feature of \BEAV.   The contents of two
windows are compared byte for byte from the current cursor
position in each window.   There must be exactly two windows to
use this command.   These windows can be displaying the same or
different buffers.   When a difference is found the cursor in
each window is moved to that position and both windows are moved
accordingly.   The display mode does not affect the operation of
this command except in restricting the cursor position to whole
units.

	\key{Esc X}{extended-command}{}

      If any command looses its binding, this command allows the
unbound command to be used.   A command can loose its binding
because the binding was assigned to another command.   When this
command is given you will be prompted for a command name.   Enter
the command name that you wish to execute, it will be executed as
if you had typed its key binding.

	\key{Esc Ctl-F}{n-way-combine}{}

      The contents of other windows can be copied sequentially
into the current window.   This is useful in combining odd-even
proms into an executable image file.   To use this command create
an empty window with a buffer file name of an empty or
nonexistent file.   Read into additional windows the files that
you want to combine.   While in the empty target window, issue
the n-way-combine command.   The data in the other windows will
be read into the current window.   The next window lower on the
screen will be read first, then the one below that, etc.

      For example; if you had two files, {\tt promlow.bin} and
{\tt promhi.bin} that you wanted to combine into a file called
{\tt prom.bin}.   First issue the file-visit command (Ctl-X Ctl-V),
enter {\tt prom.bin} at the prompt.   This file should be empty of non-
existent.   Next read {\tt promlow.bin} into a new window with the
file-visit-split command (Esc U), enter {\tt promlow.bin} at the
prompt.   Open another window for {\tt promhi.bin} with the same
command.   Go to the window containing {\tt prom.bin} (empty).   Issue
the n-way-combine command.   \BEAV\ will copy the first byte from
the window immediately below the {\tt prom.bin} window and deposit it
in the destination window buffer as well as advance the dot
position in both windows.   It will advance to the next lower
window and copy a byte from there into the destination window and
advance the dot in both windows.   This process will continue
until one of the source buffers is exhausted, or the user
terminates the command.

      The user must take care that the source buffers are in the
correct order.   They are read starting at the window immediately
below the current window.   If the target window is at the bottom
of the screen then it wraps to the top.   In this way any order
can be used and changed at will.

	\key{Esc Ctl-S}{n-way-split}{}

      This command is the mirror image of the n-way-combine.
The data in the current window is distributed among the rest of
the window buffers displayed.   The current window buffer must be
the only window buffer that contains data.   If there are two
other empty window buffers then the data will be divided two
ways.   If there are five then the data will be divided five ways

	\key{Ctl-L}{refresh-screen}{}

      The screen is reprinted from \BEAV's internal buffer.   This
is useful if the display is messed up due to transmission errors.
On a PC this is unlikely to happen.

	\key{Esc Ctl-V}{show-version}{}

      The version and date of \BEAV\ is displayed in the command
line.

	\key{Ctl-X C}{spawn-shell}{}

      A new MSDOS command shell is created.   You can return to
\BEAV\ by typing `exit'.

	\key{Ctl-U}{repeat count}{}

      This command prompts for a number to be entered.   This
causes the next command given to be repeated by that number of
times.   This command cannot have it's binding changed and cannot
be issued using the `extended-command' function.


\section{Alphabetical list of commands by name}

\begin{center}\begin{tabular}{lll}
Command                      & Key Binding     & Manual Section \\\hline
abort-cmd                    & Ctl-G           & 3.12	\\
abort-cmd                    & Ctl-X Ctl-G     & 3.12	\\
abort-cmd                    & Esc Ctl-G       & 3.12	\\
abort-cmd                    & F10             & 3.12	\\
auto-save                    & Esc A           & 3.12	\\
bind-to-key                  & Esc K           & 3.11	\\
binding-for-key              & Ctl-X ?         & 3.11	\\
binding-for-key              & Sh-F1           & 3.11	\\
bindings-load                & Ctl-X L         & 3.11	\\
buffer-set-file-name         & Ctl-F7          & 3.3	\\
buffer-set-file-name         & Ctl-X Ctl-F     & 3.3	\\
buffer-set-file-name         & Sh-F7           & 3.3	\\
buffer-set-name              & Esc Ctl-N       & 3.3	\\
buffer-size-lock             & Ctl-X Ctl-L     & 3.3	\\
buffers-display              & Ctl-F1          & 3.3	\\
buffers-display              & Ctl-X Ctl-B     & 3.3	\\
change-buffer                & Ctl-F2          & 3.3	\\
change-buffer                & Ctl-X B         & 3.3	\\
change-to-next-buffer        & Ctl-F4          & 3.3	\\
change-to-next-buffer        & Esc +           & 3.3	\\
change-to-prev-buffer        & Ctl-F5          & 3.3	\\
change-to-prev-buffer        & Esc -           & 3.3	\\
change-window-back           & Ctl-PageUp      & 3.5	\\
change-window-back           & Ctl-X P         & 3.5	\\
change-window-forw           & Ctl-PageDown    & 3.5	\\
change-window-forw           & Ctl-X N         & 3.5	\\
compare                      & Esc C           & 3.12	\\
copy-mark-to-cursor          & Esc W           & 3.6	\\
copy-mark-to-cursor          & F7              & 3.6	\\
delete-back-char             & Backspace       & 3.6	\\
delete-back-char             & Rubout          & 3.6	\\
delete-back-unit             & Esc Ctl-K       & 3.6	\\
delete-back-unit             & Esc Rubout      & 3.6	\\
delete-forw-char             & Ctl-D           & 3.6	\\
delete-forw-char             & Delete          & 3.6	\\
delete-forw-unit             & Esc D           & 3.6	\\
delete-mark-to-cursor        & Ctl-W           & 3.6	\\
delete-mark-to-cursor        & F8              & 3.6	\\
\end{tabular}\end{center}
\newpage

\begin{center}\begin{tabular}{lll}
Command                      & Key Binding     & Manual Section \\\hline
display-ascii                & Esc Ctl-A       & 3.5	\\
display-binary               & Esc Ctl-B       & 3.5	\\
display-byte-shift           & Ctl-A           & 3.5	\\
display-bytes                & Esc 1           & 3.5	\\
display-decimal              & Esc Ctl-D       & 3.5	\\
display-double-words         & Esc 4           & 3.5	\\
display-ebcdic               & Esc Ctl-E       & 3.5	\\
display-hex                  & Esc Backspace   & 3.5	\\
display-octal                & Esc Ctl-O       & 3.5	\\
display-swap-order           & Ctl-E           & 3.5	\\
display-words                & Esc 2           & 3.5	\\
extended-command             & Esc X           & 3.12	\\
file-read                    & Ctl-X Ctl-R     & 3.4	\\
file-read                    & Sh-F2           & 3.4	\\
file-save                    & Ctl-X Ctl-S     & 3.4	\\
file-save                    & Sh-F3           & 3.4	\\
file-view                    & Ctl-X V         & 3.4	\\
file-visit                   & Ctl-X Ctl-V     & 3.4	\\
file-visit                   & Sh-F4           & 3.4	\\
file-visit-split             & Esc U           & 3.4	\\
file-write                   & Ctl-X Ctl-W     & 3.4	\\
file-write                   & Sh-F5           & 3.4	\\
help                         & Esc ?           & 3.1	\\
help                         & F1              & 3.1	\\
insert-file                  & Ctl-F8          & 3.4	\\
insert-file                  & Ctl-X Tab       & 3.4	\\
insert-file                  & Sh-F8           & 3.4	\\
insert-literally             & Ctl-Q           & 3.6	\\
insert-literally             & Esc Q           & 3.6	\\
insert-toggle                & Ctl-X Ctl-A     & 3.6	\\
insert-toggle                & Insert          & 3.6	\\
insert-unit                  & Ctl-X I         & 3.6	\\
kill-buffer                  & Ctl-F3          & 3.3	\\
kill-buffer                  & Ctl-X K         & 3.3	\\
macro-end                    & Ctl-X )         & 3.10	\\
macro-execute                & Ctl-X E         & 3.10	\\
macro-start                  & Ctl-X (         & 3.10	\\
mark-set                     & Esc .           & 3.2  \\
mark-set                     & F2              & 3.2
\end{tabular}\end{center}
\newpage

\begin{center}\begin{tabular}{lll}
Command                      & Key Binding     & Manual Section \\\hline
move-back-char               & Ctl-B           & 3.2	\\
move-back-char               & West            & 3.2	\\
move-back-line               & Ctl-P           & 3.2	\\
move-back-line               & North           & 3.2	\\
move-back-page               & Esc V           & 3.2	\\
move-back-page               & PageDown        & 3.2	\\
move-back-unit               & Ctl-West        & 3.2	\\
move-back-unit               & Esc B           & 3.2	\\
move-forw-char               & Ctl-F           & 3.2	\\
move-forw-char               & East            & 3.2	\\
move-forw-line               & Ctl-N           & 3.2	\\
move-forw-line               & South           & 3.2	\\
move-forw-page               & Ctl-V           & 3.2	\\
move-forw-page               & PageUp          & 3.2	\\
move-forw-unit               & Ctl-East        & 3.2	\\
move-forw-unit               & Esc F           & 3.2	\\
move-forw-unit               & Sh-Tab          & 3.2	\\
move-to-beginning            & Esc $<$         & 3.2	\\
move-to-beginning            & Home            & 3.2	\\
move-to-buffer-split         & Esc G           & 3.2	\\
move-to-byte                 & Ctl-X G         & 3.2	\\
move-to-byte                 & F9              & 3.2	\\
move-to-end                  & End             & 3.2	\\
move-to-end                  & Esc $>$         & 3.2	\\
move-window-down             & Ctl-X Ctl-N     & 3.2	\\
move-window-down             & Ctl-Z           & 3.2	\\
move-window-up               & Ctl-X Ctl-P     & 3.2	\\
move-window-up               & Esc Z           & 3.2	\\
n-way-combine                & Esc Ctl-F       & 3.12	\\
n-way-split                  & Esc Ctl-S       & 3.12	\\
print-mark-to-cursor         & Ctl-Print       & 3.9	\\
print-mark-to-cursor         & Esc P           & 3.9	\\
quit-no-save                 & Ctl-C           & 3.8	\\
quit-no-save                 & Ctl-F10         & 3.8	\\
quit-no-save                 & Ctl-X Ctl-C     & 3.8	\\
quit-no-save                 & Sh-F10          & 3.8	\\
quit-save-all                & Ctl-F9          & 3.8	\\
quit-save-all                & Ctl-X Ctl-E     & 3.8	\\
quit-save-all                & Sh-F9           & 3.8	\\
\end{tabular}\end{center}
\newpage

\begin{center}\begin{tabular}{lll}
Command                      & Key Binding     & Manual Section \\\hline
recall-srch-string           & Ctl-R           & 3.7	\\
refresh-screen               & Ctl-L           & 3.12	\\
replace                      & Esc \%          & 3.7	\\
replace                      & F5              & 3.7	\\
save-all-buffers             & Ctl-X Return    & 3.4	\\
save-all-buffers             & Sh-F6           & 3.4	\\
save-mark-to-cursor          & Esc O           & 3.3	\\
search-again                 & Esc T           & 3.7	\\
search-again                 & F4              & 3.7	\\
search-back                  & Esc R           & 3.7	\\
search-forw                  & Esc S           & 3.7	\\
search-forw                  & F3              & 3.7	\\
show-position                & Ctl-X =         & 3.2	\\
show-save-buf                & Esc Ctl-W       & 3.3	\\
show-version                 & Esc Ctl-V       & 3.12	\\
spawn-shell                  & Ctl-X C         & 3.12	\\
swap-cursor-and-mark         & Ctl-X Ctl-X     & 3.2	\\
unit-twiddle                 & Ctl-T           & 3.6	\\
window-delete                & Ctl-X 0         & 3.5	\\
window-enlarge               & Ctl-X Z         & 3.5	\\
window-link                  & Esc L           & 3.2	\\
window-reposition            & Esc !           & 3.5	\\
window-shrink                & Ctl-X Ctl-Z     & 3.5	\\
window-single                & Ctl-X 1         & 3.5	\\
window-split                 & Ctl-X 2         & 3.5	\\
yank                         & Ctl-Y           & 3.6	\\
yank                         & F6              & 3.6	\\
yank-buffer                  & Ctl-F6          & 3.3	\\
yank-buffer                  & Esc Y           & 3.3
\end{tabular}\end{center}
\newpage


\section{Alphabetical list of commands by key binding}

\begin{center}\begin{tabular}{lll}
Command                      & Key Binding     & Manual Section \\\hline
delete-back-char             & Backspace       & 3.6	\\
display-byte-shift           & Ctl-A           & 3.5	\\
move-back-char               & Ctl-B           & 3.2	\\
quit-no-save                 & Ctl-C           & 3.8	\\
delete-forw-char             & Ctl-D           & 3.6	\\
display-swap-order           & Ctl-E           & 3.5	\\
move-forw-unit               & Ctl-East        & 3.2	\\
move-forw-char               & Ctl-F           & 3.2	\\
buffers-display              & Ctl-F1          & 3.3	\\
change-buffer                & Ctl-F2          & 3.3	\\
kill-buffer                  & Ctl-F3          & 3.3	\\
change-to-next-buffer        & Ctl-F4          & 3.3	\\
change-to-prev-buffer        & Ctl-F5          & 3.3	\\
yank-buffer                  & Ctl-F6          & 3.3	\\
buffer-set-file-name         & Ctl-F7          & 3.3	\\
insert-file                  & Ctl-F8          & 3.4	\\
quit-save-all                & Ctl-F9          & 3.8	\\
quit-no-save                 & Ctl-F10         & 3.8	\\
abort-cmd                    & Ctl-G           & 3.12	\\
refresh-screen               & Ctl-L           & 3.12	\\
move-forw-line               & Ctl-N           & 3.2	\\
move-back-line               & Ctl-P           & 3.2	\\
change-window-forw           & Ctl-PageDown    & 3.5	\\
change-window-back           & Ctl-PageUp      & 3.5	\\
print-mark-to-cursor         & Ctl-Print       & 3.9	\\
insert-literally             & Ctl-Q           & 3.6	\\
recall-srch-string           & Ctl-R           & 3.7	\\
unit-twiddle                 & Ctl-T           & 3.6	\\
move-forw-page               & Ctl-V           & 3.2	\\
delete-mark-to-cursor        & Ctl-W           & 3.6	\\
move-back-unit               & Ctl-West        & 3.2	\\
macro-start                  & Ctl-X (         & 3.10	\\
macro-end                    & Ctl-X )         & 3.10	\\
window-delete                & Ctl-X 0         & 3.5	\\
window-single                & Ctl-X 1         & 3.5	\\
window-split                 & Ctl-X 2         & 3.5	\\
\end{tabular}\end{center}
\newpage

\begin{center}\begin{tabular}{lll}
Command                      & Key Binding     & Manual Section \\\hline
show-position                & Ctl-X =         & 3.2	\\
binding-for-key              & Ctl-X ?         & 3.11	\\
change-buffer                & Ctl-X B         & 3.3	\\
spawn-shell                  & Ctl-X C         & 3.12	\\
insert-toggle                & Ctl-X Ctl-A     & 3.6	\\
buffers-display              & Ctl-X Ctl-B     & 3.3	\\
quit-no-save                 & Ctl-X Ctl-C     & 3.8	\\
quit-save-all                & Ctl-X Ctl-E     & 3.8	\\
buffer-set-file-name         & Ctl-X Ctl-F     & 3.3	\\
abort-cmd                    & Ctl-X Ctl-G     & 3.12	\\
buffer-size-lock             & Ctl-X Ctl-L     & 3.3	\\
move-window-down             & Ctl-X Ctl-N     & 3.2	\\
move-window-up               & Ctl-X Ctl-P     & 3.2	\\
file-read                    & Ctl-X Ctl-R     & 3.4	\\
file-save                    & Ctl-X Ctl-S     & 3.4	\\
file-visit                   & Ctl-X Ctl-V     & 3.4	\\
file-write                   & Ctl-X Ctl-W     & 3.4	\\
swap-cursor-and-mark         & Ctl-X Ctl-X     & 3.2	\\
window-shrink                & Ctl-X Ctl-Z     & 3.5	\\
macro-execute                & Ctl-X E         & 3.10	\\
move-to-byte                 & Ctl-X G         & 3.2	\\
insert-unit                  & Ctl-X I         & 3.6	\\
kill-buffer                  & Ctl-X K         & 3.3	\\
bindings-load                & Ctl-X L         & 3.11	\\
change-window-forw           & Ctl-X N         & 3.5	\\
change-window-back           & Ctl-X P         & 3.5	\\
save-all-buffers             & Ctl-X Return    & 3.4	\\
insert-file                  & Ctl-X Tab       & 3.4	\\
file-view                    & Ctl-X V         & 3.4	\\
window-enlarge               & Ctl-X Z         & 3.5	\\
yank                         & Ctl-Y           & 3.6	\\
move-window-down             & Ctl-Z           & 3.2	\\
delete-forw-char             & Delete          & 3.6	\\
move-forw-char               & East            & 3.2	\\
move-to-end                  & End             & 3.2	\\
window-reposition            & Esc !           & 3.5	\\
replace                      & Esc \%          & 3.7	\\
change-to-next-buffer        & Esc +           & 3.3	\\
change-to-prev-buffer        & Esc -           & 3.3	\\
\end{tabular}\end{center}
\newpage

\begin{center}\begin{tabular}{lll}
Command                      & Key Binding     & Manual Section \\\hline
mark-set                     & Esc .           & 3.2	\\
display-bytes                & Esc 1           & 3.5	\\
display-words                & Esc 2           & 3.5	\\
display-double-words         & Esc 4           & 3.5	\\
move-to-beginning            & Esc $<$         & 3.2	\\
move-to-end                  & Esc $>$         & 3.2	\\
help                         & Esc ?           & 3.1	\\
auto-save                    & Esc A           & 3.12	\\
move-back-unit               & Esc B           & 3.2	\\
display-hex                  & Esc Backspace   & 3.5	\\
Compare                      & Esc C           & 3.12	\\
display-ascii                & Esc Ctl-A       & 3.5	\\
display-binary               & Esc Ctl-B       & 3.5	\\
display-decimal              & Esc Ctl-D       & 3.5	\\
display-ebcdic               & Esc Ctl-E       & 3.5	\\
n-way-combine                & Esc Ctl-F       & 3.12	\\
abort-cmd                    & Esc Ctl-G       & 3.12	\\
delete-back-unit             & Esc Ctl-K       & 3.6	\\
buffer-set-name              & Esc Ctl-N       & 3.3	\\
display-octal                & Esc Ctl-O       & 3.5	\\
n-way-split                  & Esc Ctl-S       & 3.12	\\
show-version                 & Esc Ctl-V       & 3.12	\\
show-save-buf                & Esc Ctl-W       & 3.3	\\
delete-forw-unit             & Esc D           & 3.6	\\
move-forw-unit               & Esc F           & 3.2	\\
move-to-buffer-split         & Esc G           & 3.2	\\
bind-to-key                  & Esc K           & 3.11	\\
window-link                  & Esc L           & 3.2	\\
save-mark-to-cursor          & Esc O           & 3.3	\\
print-mark-to-cursor         & Esc P           & 3.9	\\
insert-literally             & Esc Q           & 3.6	\\
search-back                  & Esc R           & 3.7	\\
delete-back-unit             & Esc Rubout      & 3.6	\\
search-forw                  & Esc S           & 3.7	\\
search-again                 & Esc T           & 3.7	\\
file-visit-split             & Esc U           & 3.4	\\
move-back-page               & Esc V           & 3.2	\\
copy-mark-to-cursor          & Esc W           & 3.6	\\
extended-command             & Esc X           & 3.12	\\
yank-buffer                  & Esc Y           & 3.3	\\
move-window-up               & Esc Z           & 3.2	\\
\end{tabular}\end{center}
\newpage

\begin{center}\begin{tabular}{lll}
Command                      & Key Binding     & Manual Section \\\hline
help                         & F1              & 3.1	\\
mark-set                     & F2              & 3.2	\\
search-forw                  & F3              & 3.7	\\
search-again                 & F4              & 3.7	\\
replace                      & F5              & 3.7	\\
yank                         & F6              & 3.6	\\
copy-mark-to-cursor          & F7              & 3.6	\\
delete-mark-to-cursor        & F8              & 3.6	\\
move-to-byte                 & F9              & 3.2	\\
abort-cmd                    & F10             & 3.12	\\
move-to-beginning            & Home            & 3.2	\\
insert-toggle                & Insert          & 3.6	\\
move-back-line               & North           & 3.2	\\
move-back-page               & PageDown        & 3.2	\\
move-forw-page               & PageUp          & 3.2	\\
delete-back-char             & Rubout          & 3.6	\\
binding-for-key              & Sh-F1           & 3.11	\\
file-read                    & Sh-F2           & 3.4	\\
file-save                    & Sh-F3           & 3.4	\\
file-visit                   & Sh-F4           & 3.4	\\
file-write                   & Sh-F5           & 3.4	\\
save-all-buffers             & Sh-F6           & 3.4	\\
buffer-set-file-name         & Sh-F7           & 3.3	\\
insert-file                  & Sh-F8           & 3.4	\\
quit-save-all                & Sh-F9           & 3.8	\\
quit-no-save                 & Sh-F10          & 3.8	\\
move-forw-unit               & Sh-Tab          & 3.2	\\
move-forw-line               & South           & 3.2	\\
move-back-char               & West            & 3.2
\end{tabular}\end{center}


\newpage

\section{Release notes}

      Version 1.20 (3/10/91) of \BEAV\ contains the following fixes
and enhancements;

\begin{itemize}
\item       Under \UNIX\ files are created with read/write
      permissions.

\item       Fixed the bug in the terminal I/O routine that caused
      \BEAV\ to spin rather than give up control when waiting for a
      character.

\item       Added the ANSI \#define that was missing for MSDOS.

\item       Changed the D16 \#define to a unsigned short.

\item       Called ttclose on error exit.

\item       Check and limit ncol and nrow to the actual screen
      array size.

\item       Add the ability to load key bindings from a file
      automatically under MSDOS and \UNIX.

\item       Add delete current window command.

\item       Support VT100 type function keys.
\end{itemize}
      Version 1.30 (7/1/91) of \BEAV\ contains the following fixes
and enhancements;

\begin{itemize}
\item       Under MSDOS and 16 bit \UNIX\ systems the kill or copy
      region could not be over 64K bytes.   This limit has been
      eliminated.

\item       The save buffer can be made visible with the Esc Ctl-
      W command.   The save buffer is not editable.

\item       All memory allocation errors now pause and ask for
      conformation before continuing.   In previous releases only
      an error message was printed.   Since an allocation error
      generally means data loss, I have forced the user to
      respond.   Memory allocation errors are not otherwise fatal
      to \BEAV, they are probably fatal to the user's data.   The
      decision is left to the user with the appropriate warning.

\item       Two commands have been added to aid in working with
      PROM files; n-way-split (Esc Ctl-S) and n-way-combine (Esc
      Ctl-F).

\item       The speed of the delete-mark-to-cursor (Ctl-W)
      command has been greatly improved.


\item       All commands that can potentially take a lot of time
      can be stopped by pressing Ctl-G.
\end{itemize}
      Version 1.31 (11/2/91) of \BEAV\ contains the following
fixes;

\begin{itemize}
\item       A serious bug that causes a crash on systems that
      trapped the use of dereferenced pointers has been fixed.

\item       \BEAV\ now names the backup file properly under \UNIX\.
      Previously; if a dot file (.$<$filename$>$) was edited, the
      backup file was given a garbage name.   Now, a backup file
      simply has ``{\tt .bak}'' appended to the file name.

\item       You can use the buffers-display (Ctl-X, Ctl-B)
      command to; go to, kill, or save a buffer.

\item       A compile flag for DEC ULTRA was created and a
      makeable is included in this release ({\tt makefile.utx}).

\item       When a large region was deleted the offset value was
      displayed wrong, this is now fixed.

\item       A bug in the parse\_f\_name that trashed a variable is
      now fixed.

\item       Regions of never used code have been deleted.

\item       Under \UNIX\ the file permissions are maintained
      correctly when the file is saved.

\item       A number of un-niceities that lint reported have been
      fixed.
\end{itemize}
      Version 1.32 (11/8/91) of \BEAV\ contains the following
enhancements;

\begin{itemize}
\item       \BEAV\ will now compile and run on the Amiga computer.
\end{itemize}


\section{Source Availability}

      The source and MSDOS executable \BEAV\ has been posted on
news to comp.sources.misc.

      The MSDOS executable has been posted to the
\linebreak comp.binaries.ibm.pc news group.   This is archived at SIMTEL20
in \linebreak {\tt PD1:$<$MSDOS.FILUTL$>$BEAV131.ZIP}.

      If anyone does not have access to usenet, I will mail a
copy of the source on floppy for \$20.00 copying charge.   The
floppies can be in MSDOS file format or \UNIX\ tar format.   I can
also supply either QIC-24, QIC-120, QIC-150, or 9 track reel to
reel tape.   The price for the tape will include the cost of the
media.


\cleardoublepage
\pagenumbering{roman}
\tableofcontents
\cleardoublepage

\end{document}
